% !TeX program = xelatex
% Compile twice with XeLaTeX to generate the table of contents and references.
\documentclass[12pt,a4paper]{article}

\usepackage[a4paper,top=2.4cm,bottom=2.4cm,left=2.5cm,right=2.5cm]{geometry}
\usepackage{fontspec}
\usepackage{polyglossia}
\setdefaultlanguage{arabic}
\setotherlanguage{english}
\newfontfamily\arabicfont[Script=Arabic,Scale=1.08]{Amiri}
\newfontfamily\arabicfontsf[Script=Arabic,Scale=1.05]{Amiri}
\newfontfamily\englishfont{TeX Gyre Termes}

\usepackage{setspace}
\usepackage{array,tabularx,longtable,booktabs}
\usepackage{enumitem}
\usepackage{titlesec}
\usepackage{fancyhdr}
\usepackage{graphicx}
\usepackage{xcolor}
\usepackage{hyperref}
\usepackage{float}
\usepackage{listings}
\usepackage{caption}

\hypersetup{
  colorlinks=true,
  linkcolor=blue!50!black,
  urlcolor=blue!60!black,
  pdftitle={وثيقة الرؤية النهائية لمشروع نظام إدارة التدريب الميداني COOP},
  pdfauthor={مجموعة المشروع - كلية إدارة الأعمال بينبع}
}
\setstretch{1.35}
\setlength{\parskip}{5pt}
\setlength{\parindent}{0pt}
\setlist[itemize]{itemsep=2pt,topsep=3pt,rightmargin=1.2em}
\setlist[enumerate]{itemsep=3pt,topsep=3pt,rightmargin=1.2em}
\renewcommand{\arraystretch}{1.35}
\newcolumntype{Y}{>{\raggedleft\arraybackslash}X}
\newcolumntype{C}{>{\centering\arraybackslash}p{2.4cm}}

\pagestyle{fancy}
\fancyhf{}
\fancyhead[R]{نظام إدارة التدريب الميداني (COOP)}
\fancyhead[L]{وثيقة رؤية المشروع}
\fancyfoot[C]{\thepage}
\renewcommand{\headrulewidth}{0.4pt}

\titleformat{\section}{\Large\bfseries}{\thesection}{0.6em}{}
\titleformat{\subsection}{\large\bfseries}{\thesubsection}{0.6em}{}
\titleformat{\subsubsection}{\normalsize\bfseries}{\thesubsubsection}{0.6em}{}

\begin{document}

%---------------- Cover ----------------
\begin{titlepage}
\centering
\vspace*{1.5cm}
{\Large جامعة طيبة}\\[0.35cm]
{\large كلية إدارة الأعمال بينبع}\\[1.2cm]
{\Large مقرر برمجة الأعمال المتقدمة (MIS412)}\\[2.0cm]
{\Huge\bfseries وثيقة الرؤية النهائية للمشروع}\\[0.7cm]
{\LARGE\bfseries نظام إدارة التدريب الميداني}\\[0.3cm]
{\Large\textenglish{Field Training Management System (COOP)}}\\[1.8cm]
\begin{tabular}{rl}
\textbf{إعداد:} & المجموعة (ثلاثة طلاب/طالبات)\\[0.45cm]
\textbf{أعضاء المجموعة:} & 1. الاسم: \underline{\hspace{5cm}} \quad الرقم الأكاديمي: \underline{\hspace{2cm}}\\[0.45cm]
& 2. الاسم: \underline{\hspace{5cm}} \quad الرقم الأكاديمي: \underline{\hspace{2cm}}\\[0.45cm]
& 3. الاسم: \underline{\hspace{5cm}} \quad الرقم الأكاديمي: \underline{\hspace{2cm}}\\[0.45cm]
\textbf{مقدمة إلى:} & المشرف/ة على المقرر\\[0.45cm]
\textbf{التاريخ:} & \underline{\hspace{4cm}}
\end{tabular}
\vfill
{\large العام الجامعي: \underline{\hspace{3cm}}}
\end{titlepage}

\tableofcontents
\newpage

\section{معلومات المشروع}
\begin{tabularx}{\textwidth}{|p{3.5cm}|Y|}
\hline
\textbf{البند} & \textbf{التفاصيل}\\ \hline
اسم المشروع & نظام إدارة التدريب الميداني لكلية إدارة الأعمال بينبع\\ \hline
المقرر & برمجة الأعمال المتقدمة (MIS412) — المشروع (15 درجة)\\ \hline
لغة التطوير & Java باستخدام واجهات Swing\\ \hline
قاعدة البيانات & Apache Derby (Java DB) عبر JDBC\\ \hline
مصدر البيانات & ملف Excel الخاص ببيانات التدريب الميداني، بعد تهيئته للاستيراد بصيغة CSV\\ \hline
بيئة التطوير & NetBeans IDE\\ \hline
نمط التشغيل & تطبيق مكتبي محلي باستخدام Derby Embedded\\ \hline
فريق العمل & ثلاثة أعضاء\\ \hline
\end{tabularx}

\section{المقدمة ومبررات المشروع}
تُحفظ بيانات التدريب الميداني للطلاب حاليًا في ملف Excel يضم بيانات مثل اسم الطالب/ـة، والتخصص، والمرشد الأكاديمي، وجهة التدريب، والشطر عند توفره. ويؤدي الاعتماد على ملف واحد إلى صعوبة البحث والاستعلام السريع، كما يجعل تحديث البيانات واستخراج الإحصائيات عرضة للأخطاء والعمل اليدوي.

يقترح المشروع تطوير تطبيق مكتبي بواجهات رسومية وقاعدة بيانات علائقية منظمة، يتيح إدارة بيانات التدريب الميداني والاستعلام عنها بكفاءة، مع تطبيق مفاهيم البرمجة وتقنيات قواعد البيانات التي يتناولها مقرر MIS412.

\section{الرؤية والرسالة}
\subsection{الرؤية}
بناء بيئة رقمية موحدة وموثوقة لإدارة بيانات التدريب الميداني، تسهم في تحسين كفاءة العمليات الأكاديمية وتسهيل الوصول إلى المعلومات الدقيقة.

\subsection{الرسالة}
توفير نظام مكتبي بسيط وفعال يمكّن منسق التدريب من تنظيم بيانات الطلاب وتحديثها واسترجاعها، ويتيح عرض بيانات التدريب ذات الصلة للمستخدم المخوّل، مع المحافظة على سلامة البيانات وخصوصيتها.

\section{الأهداف}
يهدف المشروع إلى:
\begin{enumerate}
\item بناء قاعدة بيانات علائقية منظمة انطلاقًا من ملف البيانات الحالي.
\item تطوير واجهات رسومية واضحة لإدارة سجلات التدريب الميداني والاستعلام عنها.
\item تمكين منسق التدريب من إضافة السجلات وتعديلها وحذفها والبحث فيها وتصفيتها.
\item إتاحة عرض بيانات الطالب المتعلقة بالتدريب، وفق آلية الوصول المعتمدة في نطاق المشروع.
\item توفير إحصائيات عن توزيع الطلاب حسب جهات التدريب والتخصصات والمرشدين عند الحاجة.
\item تطبيق مفاهيم المقرر عمليًا، ومنها البرمجة الكائنية، والاستثناءات، والتعامل مع الملفات، وواجهات Swing، والاتصال بقواعد البيانات باستخدام JDBC.
\end{enumerate}

\section{المستخدمون والأدوار}
يركز الإصدار الحالي على الأدوار التالية، مع مراعاة أن التطبيق يعمل محليًا وليس خدمة إلكترونية متاحة عبر الإنترنت:

\begin{tabularx}{\textwidth}{|p{3cm}|Y|Y|}
\hline
\textbf{المستخدم} & \textbf{المهام المسموح بها} & \textbf{حدود الوصول}\\ \hline
منسق التدريب & عرض جميع السجلات، والإضافة، والتعديل، والحذف، والبحث، والتصفية، وعرض الإحصائيات، وإدارة القوائم المرجعية ضمن النطاق المعتمد. & صلاحيات الإدارة متاحة داخل التطبيق المحلي للمستخدم المسؤول عنه.\\ \hline
الطالب & الاستعلام عن بيانات تدريبه، مثل التخصص والمرشد وجهة التدريب، إذا تم اعتماد شاشة استعلام للطالب. & يقتصر العرض على سجل الطالب نفسه، ولا يُتاح له تعديل السجلات أو استعراض بيانات بقية الطلاب.\\ \hline
مدير النظام (مستقبلاً) & إدارة إعدادات التطبيق والمستخدمين والنسخ الاحتياطية عند التوسع. & ليس جزءًا إلزاميًا من الإصدار الحالي.\\ \hline
\end{tabularx}

\textbf{قرار نطاقي وأمني:} لا يتضمن الإصدار الحالي نظام مصادقة متقدمًا أو تشغيلًا شبكيًا. لذلك، إذا لم تُنفّذ آلية تحقق من هوية الطالب، فسيُقدَّم التطبيق بوصفه أداة إدارية محلية للمنسق، ولن يُدّعى أنه يوفّر وصولًا فرديًا آمنًا للطلاب. ويجب عدم توزيع قاعدة البيانات الحقيقية أو نسخها إلى أجهزة غير مخوّلة. ويمكن إضافة تسجيل الدخول والصلاحيات في إصدار لاحق.

\section{نطاق المشروع}
\subsection{داخل النطاق}
\begin{itemize}
\item استيراد البيانات الأولية من ملف Excel بعد حفظه بصيغة CSV بترميز UTF-8.
\item إدارة سجلات الطلاب: إضافة، وتعديل، وحذف، وعرض.
\item البحث بالاسم أو الرقم الأكاديمي عند توفره في مصدر البيانات.
\item التصفية حسب التخصص أو المرشد الأكاديمي أو جهة التدريب أو الشطر عند توفره.
\item إدارة القوائم المرجعية للتخصصات والمرشدين وجهات التدريب ضمن واجهة المنسق.
\item عرض لوحة تحكم وإحصائيات أساسية.
\item معالجة أخطاء الإدخال والملفات والاتصال بقاعدة البيانات برسائل واضحة.
\item حفظ البيانات محليًا واستمرارها بعد إغلاق التطبيق وإعادة تشغيله.
\end{itemize}

\subsection{خارج النطاق}
\begin{itemize}
\item تطبيقات الويب أو الهواتف الذكية.
\item الاستضافة السحابية أو الوصول عن بُعد عبر الإنترنت.
\item نظام صلاحيات ومصادقة مؤسسي متقدم في الإصدار الحالي.
\item متابعة الحضور والتقييم اليومي أو إدارة تفاصيل أداء الطالب أثناء التدريب.
\item المراسلات الفورية أو الإشعارات أو الدفع الإلكتروني.
\end{itemize}

\section{مصدر البيانات وجودتها وخصوصيتها}
تتضمن البيانات المصدرية نحو 106 سجلات وفق الملف المتاح عند إعداد هذه الوثيقة، وتشمل الشطر، واسم الطالب/ـة، والمرشد الأكاديمي، والتخصص (الإدارة، والمحاسبة، ونظم المعلومات الإدارية)، وجهة التدريب، مثل الهيئة الملكية والمستشفيات والغرفة التجارية وإدارة التعليم والبلدية والشركات الصناعية. ويُعد العدد تقريبيًا، ويُراجع عند الاستيراد الفعلي.

\subsection{قواعد تنظيف البيانات}
\begin{enumerate}
\item توحيد كتابة أسماء الجهات والتخصصات والمرشدين، بما في ذلك اختلافات الإملاء والمسافات.
\item إزالة المسافات الزائدة والتحقق من الحقول الإلزامية.
\item التحقق من السجلات المكررة، واستخدام الرقم الأكاديمي معرّفًا فريدًا إذا كان متاحًا وموثوقًا.
\item تسجيل السجلات غير الصالحة في تقرير استيراد بدل تجاهلها دون توضيح.
\item عدم اختلاق بيانات مفقودة؛ تُترك فارغة أو تُعالج وفق قرار منسق التدريب.
\end{enumerate}

\subsection{آلية الاستيراد}
يُقرأ ملف CSV بترميز UTF-8، ثم تُفحص أعمدته وقيمه قبل إدخالها إلى قاعدة البيانات. ينفّذ الاستيراد داخل معاملة قاعدة بيانات \textenglish{Transaction} قدر الإمكان لضمان اتساق البيانات. ويجب أن تتضمن العملية:
\begin{itemize}
\item التحقق من وجود الملف وصحة تنسيقه ورؤوس الأعمدة.
\item التحقق من الحقول المطلوبة والعلاقات المرجعية.
\item منع تكرار السجلات، وعدم إعادة استيراد البيانات نفسها عند كل تشغيل.
\item عرض ملخص بعد الاستيراد يتضمن عدد السجلات المقروءة والمضافة والمرفوضة وسبب الرفض.
\item إتاحة إعادة المحاولة بعد تصحيح الملف دون إنشاء سجلات مكررة.
\end{itemize}
يُنفّذ الاستيراد الأولي مرة واحدة عند تهيئة قاعدة البيانات، ولا يُعاد تلقائيًا عند كل تشغيل. ويمكن توفير خيار إداري مستقل للاستيراد اللاحق مع قواعد منع التكرار.

\subsection{الخصوصية}
تحتوي البيانات المصدرية على معلومات طلاب حقيقية؛ لذلك تُستخدم بيانات تجريبية أو مجهّلة في لقطات الشاشة والتقرير والعرض التقديمي. وتُحفظ قاعدة البيانات الأصلية محليًا، ولا تُنشر أو تُرسل إلى جهات غير مخوّلة. كما تُتجنب طباعة البيانات الشخصية في سجلات الأخطاء، ويُقيّد الوصول إلى ملفات قاعدة البيانات ونسخها الاحتياطية.

\section{المتطلبات الوظيفية}
\begin{longtable}{|C|p{12cm}|}
\hline
\textbf{الرمز} & \textbf{المتطلب}\\ \hline
F1 & تهيئة قاعدة البيانات واستيراد البيانات الأولية من CSV عند الإعداد الأول، مع منع الاستيراد المكرر.\\ \hline
F2 & إضافة طالب/ـة جديد مع إدخال البيانات المطلوبة واختيار التخصص والمرشد وجهة التدريب من قوائم منسدلة.\\ \hline
F3 & تعديل بيانات طالب/ـة موجود، مع التحقق من صحة المدخلات.\\ \hline
F4 & حذف سجل بعد عرض رسالة تأكيد، مع تحديث الواجهات والإحصائيات.\\ \hline
F5 & البحث بالاسم أو الرقم الأكاديمي إذا كان متاحًا في البيانات.\\ \hline
F6 & تصفية السجلات حسب التخصص أو المرشد أو الجهة أو الشطر عند توفره.\\ \hline
F7 & عرض جميع السجلات في جدول واضح مع إمكانية تحديث البيانات المعروضة.\\ \hline
F8 & عرض إحصائيات عدد الطلاب حسب جهة التدريب والتخصص، ويمكن إضافة المرشد.\\ \hline
F9 & إدارة بيانات التخصصات والمرشدين وجهات التدريب، بما يمنع إدخال أسماء مرجعية مكررة.\\ \hline
F10 & عرض ملخص بنتيجة استيراد CSV والأخطاء أو السجلات التي لم تُستورد.\\ \hline
F11 & تصدير نتائج البحث أو التصفية إلى CSV بوصفها ميزة إضافية إذا سمح وقت التنفيذ.\\ \hline
F12 & توفير شاشة استعلام طالب محدودة بسجله فقط إذا نُفذت آلية تحقق مناسبة؛ وإلا تُستبعد من الإصدار الحالي.\\ \hline
\end{longtable}

\section{المتطلبات غير الوظيفية}
\begin{itemize}
\item \textbf{سهولة الاستخدام:} واجهات واضحة ومتناسقة، ورسائل نجاح وتحذير وخطأ مفهومة.
\item \textbf{دعم اللغة العربية:} عرض النصوص العربية بصورة سليمة، مع ضبط اتجاه المكونات والنصوص من اليمين إلى اليسار قدر الإمكان.
\item \textbf{سلامة البيانات:} استخدام المفاتيح الأساسية والأجنبية والقيود المناسبة والتحقق من المدخلات قبل الحفظ.
\item \textbf{الأداء:} عرض مجموعة البيانات الحالية (نحو 106 سجلات) والبحث فيها بسلاسة على جهاز مناسب للمقرر.
\item \textbf{الاعتمادية:} التعامل مع فشل الاتصال أو غياب ملف الاستيراد دون انهيار التطبيق، وإظهار رسالة مفهومة.
\item \textbf{قابلية الصيانة:} فصل الواجهات عن منطق الأعمال والوصول إلى البيانات، وتسمية الأصناف والطرق بوضوح.
\item \textbf{التوافق:} توثيق إصدار Java ومكتبات Derby المطلوبة وخطوات التشغيل على جهاز آخر.
\item \textbf{الأمان:} استخدام \textenglish{PreparedStatement}، والتحقق من المدخلات، وتقييد الوصول إلى ملفات قاعدة البيانات.
\item \textbf{استمرارية البيانات:} الاحتفاظ بالتغييرات بعد إغلاق التطبيق وإعادة تشغيله.
\end{itemize}

\section{تصميم قاعدة البيانات}
تُصمم قاعدة البيانات بصورة علائقية، ويستهدف التصميم الوصول إلى الصيغة الطبيعية الثالثة (3NF) وفق طبيعة البيانات الفعلية. ويتكون النموذج الأساسي من أربعة جداول:

\begin{tabularx}{\textwidth}{|p{3cm}|Y|Y|}
\hline
\textbf{الجدول} & \textbf{الحقول الرئيسة المقترحة} & \textbf{الغرض والعلاقات}\\ \hline
\textenglish{specialization} & \textenglish{spec\_id} مفتاح أساسي، \textenglish{spec\_name} فريد وغير فارغ. & تخزين التخصصات، ويرتبط التخصص الواحد بعدة طلاب.\\ \hline
\textenglish{advisor} & \textenglish{advisor\_id} مفتاح أساسي، \textenglish{advisor\_name} فريد وغير فارغ. & تخزين المرشدين الأكاديميين، ويرتبط المرشد الواحد بعدة طلاب.\\ \hline
\textenglish{organization} & \textenglish{org\_id} مفتاح أساسي، \textenglish{org\_name} فريد وغير فارغ. & تخزين جهات التدريب، وترتبط الجهة الواحدة بعدة طلاب.\\ \hline
\textenglish{student} & \textenglish{student\_id} مفتاح أساسي، \textenglish{student\_number} رقم أكاديمي فريد عند توفره، \textenglish{full\_name}، \textenglish{section}، \textenglish{gender} عند توفرهما، وحقول المفاتيح الأجنبية. & تخزين سجل الطالب وربطه بالتخصص والمرشد وجهة التدريب.\\ \hline
\end{tabularx}

\subsection{قواعد التصميم والعلاقات}
\begin{itemize}
\item علاقة واحد إلى متعدد من \textenglish{specialization} إلى \textenglish{student}.
\item علاقة واحد إلى متعدد من \textenglish{advisor} إلى \textenglish{student}.
\item علاقة واحد إلى متعدد من \textenglish{organization} إلى \textenglish{student}.
\item يكون \textenglish{student\_number} فريدًا إذا كان الرقم الأكاديمي موجودًا في المصدر، ولا يُستخدم اسم الطالب وحده مفتاحًا فريدًا.
\item يُفضّل منع حذف سجل مرجعي مرتبط بطلاب، أو مطالبة المستخدم بمعالجة الارتباطات أولًا؛ ولا تُحذف سجلات الطلاب تلقائيًا بسبب حذف جهة أو مرشد.
\item تُستخدم قيود \textenglish{NOT NULL} و\textenglish{UNIQUE} والمفاتيح الأجنبية بحسب الحقول المتاحة فعليًا.
\end{itemize}

يُرفق بالتقرير مخطط الكيانات والعلاقات \textenglish{ER Diagram} وسكربت إنشاء الجداول والقيود. ويُراجع توصيف الحقول النهائي بعد فحص ملف Excel الأصلي؛ فلا يُضاف حقل غير موجود أو غير مطلوب دون مبرر.

\subsection{الاتصال بقاعدة البيانات}
يُستخدم Apache Derby في الوضع المضمّن \textenglish{Embedded} لتسهيل تشغيل التطبيق وتسليمه محليًا. ويمكن دراسة وضع الخادم \textenglish{Network Server} إذا تطلبت بيئة التنفيذ ذلك، باستخدام عنوان اتصال مثل:
\begin{center}
\textenglish{\small jdbc:derby://localhost:1527/CoopDB}
\end{center}
ولا يُعد وضع الخادم جزءًا إلزاميًا من الإصدار الحالي.

\section{البنية البرمجية المقترحة}
تعتمد البنية على فصل المسؤوليات لتسهيل التطوير والاختبار والصيانة:

\begin{tabularx}{\textwidth}{|p{2.7cm}|Y|Y|}
\hline
\textbf{الحزمة} & \textbf{المحتوى} & \textbf{المسؤولية}\\ \hline
\textenglish{model} & \textenglish{Person, Student, Advisor, Organization, Specialization} & تمثيل الكيانات والبيانات. يرث \textenglish{Student} و\textenglish{Advisor} من \textenglish{Person}، بينما تكون \textenglish{Organization} و\textenglish{Specialization} كيانين مستقلين.\\ \hline
\textenglish{dao} & \textenglish{DBConnection, StudentDAO, AdvisorDAO, OrganizationDAO, SpecializationDAO} & تنفيذ عمليات قاعدة البيانات باستخدام JDBC.\\ \hline
\textenglish{service} & \textenglish{StudentService, ImportService, StatisticsService} & منطق الأعمال والتحقق من قواعد النظام والتنسيق بين الواجهة وDAO.\\ \hline
\textenglish{exception} & \textenglish{DataAccessException, ValidationException, ImportException} & تعريف الاستثناءات الخاصة ومعالجة الأخطاء.\\ \hline
\textenglish{util} & \textenglish{CSVReader, DataLoader, ValidationUtil} & أدوات قراءة الملفات والتحقق والتهيئة.\\ \hline
\textenglish{ui} & النوافذ والشاشات ومكونات Swing & عرض البيانات واستقبال تفاعلات المستخدم.\\ \hline
\end{tabularx}

تُنفّذ عمليات قاعدة البيانات في طبقة DAO، وتُطبق قواعد العمل في طبقة Service، بينما تقتصر الواجهة على عرض النتائج واستقبال المدخلات. ويُتجنب وضع استعلامات SQL أو منطق الأعمال مباشرة داخل مستمعات أحداث Swing.

\subsection{مبادئ التنفيذ}
\begin{itemize}
\item استخدام \textenglish{PreparedStatement} بدل دمج مدخلات المستخدم مباشرة في نص SQL.
\item إغلاق موارد JDBC بصورة آمنة باستخدام \textenglish{try-with-resources}.
\item استخدام المعاملات عند العمليات التي تتطلب تنفيذ عدة أوامر مترابطة.
\item تشغيل واجهة Swing عبر \textenglish{SwingUtilities.invokeLater}.
\item تنفيذ العمليات التي قد تستغرق وقتًا خارج مسار واجهة المستخدم، مثل الاستيراد أو تحميل البيانات، باستخدام آلية مناسبة مثل \textenglish{SwingWorker} لتجنب تجمد الواجهة.
\item تطبيق التحقق من المدخلات قبل إرسالها إلى قاعدة البيانات، مع معالجة الاستثناءات برسائل مناسبة.
\end{itemize}

\section{تصميم الواجهات وتجربة المستخدم}
تكون الواجهات باللغة العربية قدر الإمكان، مع مراعاة الاتساق البصري وسهولة التنقل. وتشمل الشاشات المقترحة:

\begin{tabularx}{\textwidth}{|p{3cm}|Y|Y|}
\hline
\textbf{الشاشة} & \textbf{الوظيفة} & \textbf{المكونات الرئيسة}\\ \hline
الرئيسية / لوحة التحكم & عرض ملخص سريع والتنقل إلى الوظائف. & بطاقات إجمالي الطلاب والتخصصات والجهات، وأزرار الانتقال.\\ \hline
إدارة الطلاب & إضافة السجلات وتعديلها وحذفها. & حقول إدخال، وقوائم منسدلة، وأزرار حفظ وتعديل وحذف.\\ \hline
عرض وبحث الطلاب & عرض السجلات والبحث والتصفية. & \textenglish{JTable, JScrollPane, JTextField, JComboBox}.\\ \hline
إدارة البيانات المرجعية & إدارة التخصصات والمرشدين والجهات. & جداول وقوائم ونماذج إضافة وتعديل.\\ \hline
الإحصائيات & عرض توزيع الطلاب حسب الجهة والتخصص. & مخططات أعمدة أو رسوم بسيطة باستخدام \textenglish{Graphics2D}.\\ \hline
استيراد البيانات & اختيار CSV والتحقق منه وعرض نتيجة الاستيراد. & \textenglish{JFileChooser}، ورسائل التحقق وملخص النتائج.\\ \hline
استعلام الطالب (اختياري) & عرض بيانات الطالب نفسه فقط عند وجود تحقق مناسب. & حقل تحقق وآلية عرض محدودة، وإلا تؤجل هذه الشاشة.\\ \hline
\end{tabularx}

تعمل الأزرار عبر مستمعي الأحداث \textenglish{ActionListener}، وتُراعى قابلية قراءة النصوص العربية واتجاه الواجهة. ويُحدّث الجدول والإحصائيات بعد عمليات الإضافة أو التعديل أو الحذف الناجحة.

\section{خطة العمل وتوزيع المهام}
تُنفذ الخطة على مراحل، مع إمكانية تداخل بعض المهام:

\begin{longtable}{|C|p{7.4cm}|p{2.2cm}|p{2cm}|}
\hline
\textbf{المرحلة} & \textbf{المهمة} & \textbf{المسؤول} & \textbf{المدة}\\ \hline
1 & تحليل المتطلبات وفحص ملف Excel وتنظيف البيانات وتحديد الحقول. & الجميع & الأسبوع 1\\ \hline
2 & تصميم قاعدة البيانات وإنشاء الجداول والعلاقات واستيراد البيانات الأولية. & العضو 1 & الأسبوع 1--2\\ \hline
3 & بناء الأصناف وطبقة DAO وService والاستثناءات. & العضو 2 & الأسبوع 2--3\\ \hline
4 & بناء واجهات Swing وربطها بطبقات النظام. & العضو 3 & الأسبوع 2--4\\ \hline
5 & دمج المكونات واختبار العمليات والاستيراد والإحصائيات. & الجميع & الأسبوع 4\\ \hline
6 & معالجة الملاحظات وتجهيز التقرير ولقطات الشاشة وملفات التسليم. & الجميع & الأسبوع 5\\ \hline
7 & الاستعداد للمناقشة وتنفيذ عرض عملي بمشاركة جميع الأعضاء. & الجميع & الأسبوع 5\\ \hline
8 & استكمال تقييم أعضاء المجموعة وفق متطلبات المقرر. & الجميع & بعد المناقشة\\ \hline
\end{longtable}

يتبادل أعضاء المجموعة مراجعة الكود، ويحرص كل عضو على فهم المكونات الأساسية للنظام، بما يضمن القدرة على شرح المشروع والمشاركة في المناقشة.

\section{خطة الاختبار}
تُنفذ الاختبارات على بيانات تجريبية قبل استخدام البيانات الفعلية:

\begin{longtable}{|p{6cm}|p{8cm}|}
\hline
\textbf{سيناريو الاختبار} & \textbf{النتيجة المتوقعة}\\ \hline
إضافة سجل ببيانات صحيحة & حفظ السجل وظهور رسالة نجاح وتحديث الجدول.\\ \hline
إضافة سجل مع ترك حقل إلزامي فارغًا & رفض الحفظ وإظهار رسالة توضح الحقل المطلوب دون توقف البرنامج.\\ \hline
إضافة رقم أكاديمي مستخدم مسبقًا & منع التكرار وإظهار رسالة مناسبة.\\ \hline
تعديل بيانات طالب & حفظ التغييرات وعرض القيم الجديدة.\\ \hline
حذف سجل طالب & طلب التأكيد ثم الحذف وتحديث الجدول والإحصائيات.\\ \hline
محاولة حذف جهة مرتبطة بطلاب & منع الحذف أو طلب معالجة الارتباطات أولًا، دون فقد بيانات الطلاب.\\ \hline
البحث عن اسم عربي موجود & عرض السجلات المطابقة بصورة صحيحة.\\ \hline
البحث عن قيمة غير موجودة & عرض رسالة تفيد بعدم وجود نتائج.\\ \hline
استيراد ملف CSV صحيح & استيراد السجلات الصحيحة وعرض ملخص بعدد السجلات.\\ \hline
استيراد ملف فارغ أو بتنسيق غير صحيح & رفض الملف وإظهار سبب واضح.\\ \hline
استيراد ملف يتضمن سجلات مكررة أو ناقصة & منع التكرار وتوضيح السجلات التي لم تُستورد وأسبابها.\\ \hline
إعادة تشغيل التطبيق بعد الاستيراد & بقاء البيانات دون تكرار الاستيراد الأولي.\\ \hline
إدخال نص عربي & حفظ النص وعرضه بصورة سليمة.\\ \hline
فشل الاتصال بقاعدة البيانات أو غيابها & إظهار رسالة خطأ واضحة ومعالجة الاستثناء دون انهيار غير متحكم به.\\ \hline
تحديث الإحصائيات بعد إضافة أو حذف طالب & انعكاس التغيير في الأعداد والرسوم المعروضة.\\ \hline
\end{longtable}

\section{المخاطر وخطط المعالجة}
\begin{longtable}{|p{5cm}|p{9cm}|}
\hline
\textbf{الخطر المحتمل} & \textbf{خطة المعالجة}\\ \hline
اختلاف صيغ أسماء الجهات أو المرشدين & توحيد القيم قبل الاستيراد واستخدام جداول مرجعية ذات أسماء فريدة.\\ \hline
تكرار سجلات الطلاب & استخدام الرقم الأكاديمي معرّفًا فريدًا عند توفره، والتحقق قبل الإدخال.\\ \hline
ظهور الحروف العربية بصورة غير صحيحة & استخدام UTF-8 عند قراءة CSV، واختيار إعدادات خطوط مناسبة، واختبار الإدخال والعرض العربي.\\ \hline
صعوبة تشغيل Derby على جهاز آخر & تضمين المكتبات المطلوبة وتوثيق إصدار Java وخطوات التشغيل ومسار قاعدة البيانات.\\ \hline
تجمد واجهة Swing أثناء العمليات & تنفيذ العمليات الطويلة خارج مسار واجهة المستخدم وتحديث الواجهة بعد اكتمالها.\\ \hline
حذف بيانات مرجعية مرتبطة بسجلات طلاب & تطبيق قيود العلاقات ومنع الحذف غير الآمن أو طلب معالجة الارتباطات أولًا.\\ \hline
تسريب بيانات الطلاب & استخدام بيانات مجهّلة في التقرير والعرض، وتقييد الوصول إلى قاعدة البيانات وعدم نشرها.\\ \hline
تأخر أحد أعضاء المجموعة أو ضيق الوقت & متابعة أسبوعية وتوزيع واضح للمهام ومراجعة متبادلة للكود.\\ \hline
\end{longtable}

\section{المخرجات ومعايير التقييم}
تُجهز المخرجات بما يتوافق مع توزيع الدرجات الوارد في وصف المشروع:

\begin{tabularx}{\textwidth}{|p{4cm}|C|Y|}
\hline
\textbf{معيار التقييم} & \textbf{الدرجة} & \textbf{المخرجات المتوقعة}\\ \hline
إنشاء قاعدة البيانات & 4 & سكربت إنشاء الجداول والقيود، ومخطط ER، وملف البيانات بعد تنظيفه، ودليل الاستيراد.\\ \hline
كتابة الكود & 5 & ملفات Java المصدرية، وتنظيم الحزم، والاتصال بقاعدة البيانات، وتنفيذ الوظائف المطلوبة.\\ \hline
التقرير & 3 & تقرير منسق يتضمن المقدمة والتحليل والتصميم والتنفيذ والاختبار والنتائج ولقطات الشاشة.\\ \hline
المناقشة & 3 & عرض عملي للنظام وشرح مكوناته وتنفيذ سيناريوهات اختبار والإجابة عن الأسئلة بمشاركة أعضاء المجموعة.\\ \hline
\textbf{المجموع} & \textbf{15} & \\ \hline
\end{tabularx}

\subsection{ملفات التسليم المقترحة}
\begin{itemize}
\item المشروع البرمجي كاملًا وملفاته المصدرية.
\item سكربت إنشاء قاعدة البيانات.
\item ملف CSV تجريبي خالٍ من البيانات الشخصية الحقيقية.
\item مخطط الكيانات والعلاقات.
\item التقرير النهائي بصيغة PDF.
\item ملف تعليمات التشغيل والمتطلبات البرمجية.
\item صور توضح الشاشات الأساسية ونتائج الاختبارات باستخدام بيانات تجريبية.
\end{itemize}

\section{النتائج المتوقعة}
من المتوقع أن ينتج عن المشروع تطبيق مكتبي يعمل محليًا لإدارة بيانات التدريب الميداني، ويتيح للمنسق تنظيم سجلات الطلاب وتحديثها والبحث فيها واستخراج إحصائيات أساسية، مع قاعدة بيانات علائقية منظمة وواجهات رسومية سهلة الاستخدام. كما يُتوقع أن يقدم المشروع تطبيقًا عمليًا لمفاهيم مقرر MIS412، وأن يكون قابلًا للتشغيل والعرض والمناقشة وفق نطاقه المحدد.

\section{الافتراضات والقرارات التي تتطلب تأكيدًا}
تعتمد هذه الوثيقة على المعلومات المتاحة عند إعدادها، وتبقى النقاط التالية بحاجة إلى مراجعة مع المشرف/ة أو منسق التدريب قبل التنفيذ:
\begin{enumerate}
\item التأكد من الحقول الفعلية في ملف Excel، وبخاصة وجود الرقم الأكاديمي والجنس والشطر.
\item تحديد ما إذا كان المطلوب تطبيقًا إداريًا محليًا للمنسق فقط، أم أن شاشة استعلام الطالب مطلوبة ضمن التقييم.
\item إذا طُلب وصول الطلاب إلى سجلاتهم، فيجب اعتماد آلية تحقق مناسبة قبل تنفيذ هذه الوظيفة.
\item تأكيد إصدار Java ونسخة Derby وآلية التسليم المعتمدة في المقرر.
\item تأكيد ما إذا كان تصدير النتائج إلى CSV ميزة مطلوبة أم إضافية.
\end{enumerate}

\section{الموافقة}
\begin{tabularx}{\textwidth}{|Y|p{2.5cm}|p{2.5cm}|p{2.5cm}|}
\hline
\textbf{الاسم} & \textbf{الصفة} & \textbf{التوقيع} & \textbf{التاريخ}\\ \hline
 & المشرف/ة على المقرر & & \\[0.7cm] \hline
 & ممثل المجموعة & & \\[0.7cm] \hline
\end{tabularx}

\vspace{0.5cm}
\textbf{قرار المشرف/ة:} \quad $\square$ موافقة \quad $\square$ موافقة مع ملاحظات \quad $\square$ إعادة النظر

\vspace{0.5cm}
\textbf{الملاحظات:}

\rule{\textwidth}{0.4pt}\\[0.7cm]
\rule{\textwidth}{0.4pt}\\[0.7cm]
\rule{\textwidth}{0.4pt}

\end{document}
